\documentclass[11pt]{article}
\usepackage[letterpaper,margin=0.68in,headheight=15pt]{geometry}
\usepackage[T1]{fontenc}
\usepackage{lmodern}
\usepackage{amsmath,amssymb}
\usepackage{array,booktabs,enumitem,multicol}
\usepackage{xcolor}
\usepackage{tikz}
\usepackage{fancyhdr}
\usepackage{microtype}
\setlength{\parindent}{0pt}
\setlength{\parskip}{5pt}
\definecolor{olive}{HTML}{555B35}
\definecolor{gold}{HTML}{B18A22}
\definecolor{soft}{HTML}{F4F1DE}
\definecolor{bluegray}{HTML}{E8EEF2}
\pagestyle{fancy}
\fancyhf{}
\fancyhead[L]{\textcolor{olive}{\small MATH\& 141: Precalculus I}}
\fancyhead[R]{\textcolor{gold}{\small Fall 2026}}
\fancyfoot[C]{\thepage}
\renewcommand{\headrulewidth}{0.45pt}
\renewcommand{\footrulewidth}{0pt}
\setlist[enumerate]{leftmargin=*,itemsep=0.08in,topsep=0.05in}
\newcommand{\assessmenttitle}[2]{%
  {\Large\bfseries Module #1: #2}\par\vspace{3pt}
}
\newcommand{\studentline}{%
\noindent\textbf{Name:}\hspace{2.6in}\textbf{Date:}\hspace{1.15in}\par\vspace{7pt}}
\newcommand{\directions}[1]{\textit{#1}\par\vspace{4pt}}
\newcommand{\answerblank}[1]{\vspace{#1}}
\newcommand{\blankgrid}[4]{%
\begin{center}
\begin{tikzpicture}[x=0.75cm,y=0.75cm]
  \draw[step=1,gray!28,very thin] (#1,#3) grid (#2,#4);
  \draw[->,thick] (#1,0)--(#2,0) node[right] {$x$};
  \draw[->,thick] (0,#3)--(0,#4) node[above] {$y$};
  \foreach \x in {#1,...,#2}{\ifnum\x=0\else\draw (\x,0.10)--(\x,-0.10) node[below=2pt,font=\scriptsize]{\x};\fi}
  \foreach \y in {#3,...,#4}{\ifnum\y=0\else\draw (0.10,\y)--(-0.10,\y) node[left=2pt,font=\scriptsize]{\y};\fi}
\end{tikzpicture}
\end{center}}
\newcommand{\smallgrid}[4]{%
\begin{tikzpicture}[x=0.50cm,y=0.50cm]
  \draw[step=1,gray!28,very thin] (#1,#3) grid (#2,#4);
  \draw[->,thick] (#1,0)--(#2,0) node[right] {$x$};
  \draw[->,thick] (0,#3)--(0,#4) node[above] {$y$};
\end{tikzpicture}}
\begin{document}

% MODULE 10 PREQUIZ
\assessmenttitle{10}{Matrices and Gaussian Elimination -- Prequiz}
\studentline
\directions{Four short questions. Translate freely between systems, augmented matrices, and row operations.}
\begin{enumerate}
\item Write the augmented matrix for
\[
\begin{cases}
x+2y=5,\\
3x-y=4.
\end{cases}
\]
\answerblank{0.55in}
\item Starting with rows $[1\ \ 2\mid5]$ and $[3\ \ -1\mid4]$, apply $R_2\leftarrow R_2-3R_1$. Write the new second row.
\answerblank{0.52in}
\item What does an augmented row $[0\ \ 0\mid3]$ tell you about the system?
\answerblank{0.45in}
\item In a system with more variables than pivot rows, what can a zero row $[0\ \ 0\ \ 0\mid0]$ indicate about the solution set?
\answerblank{0.52in}
\end{enumerate}
\newpage

% MODULE 10 CHECKIN PAGE 1
\assessmenttitle{10}{Matrices and Gaussian Elimination -- Weekly Check-In}
\studentline
\directions{Three questions. Translate freely between systems, augmented matrices, and row operations.}
\textbf{1. Gaussian elimination.}

Consider
\[
\begin{cases}
x+y-z=2,\\
2x-y+z=4,\\
3x+2y+z=9.
\end{cases}
\]
\begin{enumerate}[label=(\alph*),leftmargin=0.35in]
\item Write the augmented matrix.
\answerblank{0.55in}
\item Apply $R_2\leftarrow R_2-2R_1$ and $R_3\leftarrow R_3-3R_1$, and write the resulting matrix.
\answerblank{0.80in}
\item Continue using elementary row operations to solve the system. State the final ordered triple.
\answerblank{1.75in}
\end{enumerate}
\newpage

% MODULE 10 CHECKIN PAGE 2
\assessmenttitle{10}{Matrices and Gaussian Elimination -- Weekly Check-In, continued}
\studentline
\textbf{2. Classifying echelon forms.}

For each augmented matrix, classify the system as having a unique solution, no solution, or infinitely many solutions. Briefly explain what feature of the matrix determines the classification.
\[
A=\left[\begin{array}{cc|c}1&2&4\\0&1&3\end{array}\right],\qquad
B=\left[\begin{array}{cc|c}1&-1&2\\0&0&5\end{array}\right],
\]
\[
C=\left[\begin{array}{ccc|c}1&2&-1&4\\0&1&3&2\\0&0&0&0\end{array}\right].
\]
\answerblank{1.05in}

\textbf{3. Quadratic fitting by a matrix.}

A quadratic $y=ax^2+bx+c$ passes through $(0,1)$, $(1,4)$, and $(2,9)$.
\begin{enumerate}[label=(\alph*),leftmargin=0.35in]
\item Write the three linear equations for $a,b,c$ and the corresponding augmented matrix.
\answerblank{0.70in}
\item Solve for $a,b,c$ and write the quadratic function.
\answerblank{0.78in}
\end{enumerate}

\end{document}
